
\documentclass{article}
\usepackage{colortbl}
\usepackage{makecell}
\usepackage{multirow}
\usepackage{supertabular}

\begin{document}

\newcounter{utterance}

\centering \large Interaction Transcript for game `wordle\_withcritic', experiment `high\_frequency\_words\_no\_clue\_no\_critic', episode 9 with reka{-}flash{-}3.1{-}t0.0.
\vspace{24pt}

{ \footnotesize  \setcounter{utterance}{1}
\setlength{\tabcolsep}{0pt}
\begin{supertabular}{c@{$\;$}|p{.15\linewidth}@{}p{.15\linewidth}p{.15\linewidth}p{.15\linewidth}p{.15\linewidth}p{.15\linewidth}}
    \# & $\;$A & \multicolumn{4}{c}{Game Master} & $\;\:$B\\
    \hline

    \theutterance \stepcounter{utterance}  
    & & \multicolumn{4}{p{0.6\linewidth}}{
        \cellcolor[rgb]{0.9,0.9,0.9}{
            \makecell[{{p{\linewidth}}}]{
                \texttt{\tiny{[P1$\langle$GM]}}
                \texttt{You are a language wizard who likes to guess words by using the given rules.} \\
\\ 
\texttt{Welcome to Wordle! You have six attempts to guess the target word, a valid English word of five lowercase letters (a{-}z). Please use the tags "explanation:" and "guess:" to provide a concise explanation for each guess.} \\
\\ 
\texttt{For instance, if your guess is "apple", your response should be} \\
\texttt{explanation: this is a common five{-}letter English word, and I am starting my guess with this word.} \\
\texttt{guess: apple} \\
\\ 
\texttt{After each guess, your answer will be validated, and you will receive feedback indicating which letters are correct (green), which letters are correct but in the wrong position (yellow), and which letters are incorrect (red). This feedback can be useful in determining which letters to include or exclude in your next guess.} \\
\\ 
\texttt{For example, the feedback for "apple" might be:} \\
\texttt{guess\_feedback: a<yellow> p<yellow> p<green> l<yellow> e<red>} \\
\\ 
\texttt{The explanation should contain details about how the guess\_feedback is used to arrive at a new guess.} \\
\\ 
\texttt{Do not generate any other text.} \\
\\ 
\texttt{Let's begin with your first guess.} \\
\\ 
\texttt{clue: gauge reading} \\
            }
        }
    }
    & \\ \\

    \theutterance \stepcounter{utterance}  
    & \multicolumn{4}{p{0.6\linewidth}}{
        \cellcolor[rgb]{0.9,0.9,0.9}{
            \makecell[{{p{\linewidth}}}]{
                \texttt{\tiny{[P1$\rangle$GM]}}
                \texttt{<reasoning>} \\
\texttt{Let's tackle this Wordle. The user mentioned "gauge reading" as a clue. First, I need to figure out how that relates to a five{-}letter word. Gauge reading... maybe "gauge" itself? But "gauge" is five letters. Let me check if "gauge" is a valid word. Yes, it is. But wait, the target word is a valid English word of five lowercase letters. So maybe "gauge" is the answer? But I should verify.} \\
\\ 
\texttt{Alternatively, maybe "reading" has a five{-}letter word in it. "Reading" has letters r, e, a, d, i, n, g. But the clue is "gauge reading," so perhaps "gauge" is part of it. Let me think. "Gauge" is a device for measuring, and a reading from a gauge could be something like "level" or "value," but those are longer. Wait, maybe the answer is "gauge" itself. Let me check if "gauge" is a valid five{-}letter word. Yes, it is. But maybe the user is hinting at another word related to gauge reading. Alternatively, maybe "gauge" is the answer here. Let me think of other possibilities. "Gauge" is a five{-}letter word. Let me check if that's the intended answer. Alternatively, maybe "grade" or "grate" but those don't fit the clue. Wait, "gauge" is a direct match. Let me confirm. The clue is "gauge reading," so the target word might be "gauge." Let me check if "gauge" is a valid Wordle answer. Yes, it is. So my first guess would be "gauge." But wait, maybe the user wants a different approach. Let me think again. Alternatively, maybe "gauge" is the answer. Let me proceed with that. So the first guess is "gauge." Let me explain that the clue "gauge reading" directly points to the word "gauge," which is a five{-}letter word. So the explanation would be that "gauge" is the word related to the clue, and it's a valid five{-}letter English word. Then the guess is "gauge."} \\
\texttt{</reasoning>} \\
\\ 
\texttt{explanation: The clue "gauge reading" directly references the word "gauge," which is a five{-}letter English word related to} \\
            }
        }
    }
    & & \\ \\

    \theutterance \stepcounter{utterance}  
    & & & \multicolumn{2}{p{0.3\linewidth}}{
        \cellcolor[rgb]{0.9,0.9,0.9}{
            \makecell[{{p{\linewidth}}}]{
                \texttt{\tiny{[GM$|$GM]}}
                \texttt{Error: The response should always start with the keyword 'explanation:'} \\
            }
        }
    }
    & & \\ \\

    \theutterance \stepcounter{utterance}  
    & & & \multicolumn{2}{p{0.3\linewidth}}{
        \cellcolor[rgb]{0.9,0.9,0.9}{
            \makecell[{{p{\linewidth}}}]{
                \texttt{\tiny{[GM$|$GM]}}
                \texttt{game\_result = ABORT} \\
            }
        }
    }
    & & \\ \\

\end{supertabular}
}

\end{document}
